\section{离散化与变分：连续方程怎么变成 $Au=b$}

这一节全部是 \textbf{未建}（只在 \texttt{NOW.md} 规划）。它的价值不在「实现细节」，
而在\textbf{把一道 PDE 拆成五步数学动作}——面试里能说清这五步的人远少于能背公式的人。

\begin{center}
\begin{tabular}{L{1.1cm}L{3.2cm}L{4.0cm}L{5.2cm}}
\toprule
步 & 数学动作 & 产物 & 为什么不能省 \\
\midrule
1 & 选网格、定自由度 & 节点集 $\{x_i\}$、编号表 & 编号决定矩阵的稀疏结构与带宽 \\
2 & 弱形式 + 分部积分 & $\int\nabla u\cdot\nabla v=\int fv$ & 降阶：二阶导挪到测试函数上 \\
3 & 单元刚度装配 & $K_e\to A$ & 每个单元的积分只算一次 \\
4 & 载荷 $b_i=\int f\varphi_i$ & $b$ & 用闭式积分才不引入额外误差 \\
5 & 边界条件 & 消元后的 $A,b$ & 不改边界，矩阵奇异（常值在零空间）\\
\bottomrule
\end{tabular}
\end{center}

\begin{know}{F1}{结构化网格与自由度编号}{未建：math/grid.hpp}
\field{是什么} 1D：$[a,b]$ 上 $n$ 个节点、$n-1$ 个单元，$h=(b-a)/(n-1)$。
2D：$n_x\times n_y$ 个节点、$(n_x-1)(n_y-1)$ 个矩形单元，每格切两个三角形，
节点按行主序编号 $p=j\,n_x+i$。

\field{为什么需要} 编号不是记账，是\textbf{矩阵结构}：行主序编号给出带宽 $O(n_x)$ 的带状矩阵，
换成对角编号带宽变 $O(n_x+n_y)$——同一个 PDE、不同的解算代价。

\field{关键结论} 自由度 = 节点（P1 元）；边界自由度是四条边上的节点。成立条件：矩形域、等距、固定对角线方向。

\field{判据} 手算表对拍：1D $n=5,[0,1]\Rightarrow$ 节点 $\{0,.25,.5,.75,1\}$、边界 $\{0,4\}$；
2D $4\times3\Rightarrow$ 12 节点、内点 2、边界 10、三角形 $2\cdot3\cdot2=12$。

\field{前置} L1。

\field{常见误解} 以为「单元数等于节点数」。1D 单元是 $n-1$、2D 单元是 $2(n_x-1)(n_y-1)$；
把单元当节点会导致刚度矩阵维度错。
\end{know}

\begin{know}{F2}{P1 单元刚度矩阵}{未建：math/assemble.hpp}
\field{是什么} 1D：$K_e=\frac1h\begin{bmatrix}1&-1\\-1&1\end{bmatrix}$；
2D 直角三角形：$K_e=\frac12\begin{bmatrix}2&-1&-1\\-1&1&0\\-1&0&1\end{bmatrix}$。

\field{为什么需要} 它是「$-\Delta$ 在分片线性空间上的矩阵表示」，整个有限元就是从它装配出来的。

\field{关键结论} 两个必查性质：\textbf{行和为 0}（常函数在核里，对应「常数不影响梯度」）、
\textbf{对称}（因为双线性形式对称）。成立条件：P1（分片线性）基函数；单元几何规则。

\field{判据} 1D $n=3,h=0.5$：$A=\frac1h\begin{bmatrix}1&-1&0\\-1&2&-1\\0&-1&1\end{bmatrix}$ 逐项对拍；
$A(i,j)=A(j,i)$ 要求浮点\textbf{精确}相等（同一条装配路径算两次）。

\field{前置} F1。

\field{常见误解} 以为「2D 的 $K_e$ 是 1D 的张量积」。直角三角形的 $K_e$ 不是矩形元的张量积形式；
用错会得到与解析解不符的收敛阶。
\end{know}

\begin{know}{F3}{弱形式与分部积分}{未建：math/assemble.hpp}
\field{是什么} 对 $-\Delta u=f$（Dirichlet 0），两边乘测试函数 $v$ 并积分，用分部积分：
\[
\int_\Omega\nabla u\cdot\nabla v\,dx=\int_\Omega fv\,dx\quad\forall v\in H^{1}_{0}.
\]

\field{为什么需要} 强形式要求 $u$ 二阶可导；弱形式只要求一阶可积——这才让分片线性函数成为合法解。

\field{关键结论} 分部积分把「一个二阶导」换成「两个一阶导的乘积」，同时\textbf{把边界项丢给自然边界条件}。
Dirichlet 条件是本质边界（要强加），Neumann 是自然边界（出现在右端）。
成立条件：$u,v\in H^{1}$；边界足够正则（分段光滑即可）。

\field{判据} 无（这一步是推导，不是计算）；但「行和为 0」（F2）就是弱形式在常数上的推论。

\field{前置} F2。

\field{常见误解} 以为「弱形式是强形式的近似」。两者对光滑解\textbf{完全等价}；弱形式只是把要求放宽到
更差的函数类，从而让离散解有定义。
\end{know}

\begin{know}{F4}{一致载荷与闭式积分}{未建：math/assemble.hpp}
\field{是什么} $b_i=\int f\varphi_i$。1D 线性 $f$：单元上 $\frac h6[2f_i+f_{i+1},\ f_i+2f_{i+1}]$；
2D 三角形用 $\int\lambda_1^{a}\lambda_2^{b}\lambda_3^{c}\,dA=2A\frac{a!\,b!\,c!}{(a+b+c+2)!}$。

\field{为什么需要} 载荷的误差会\textbf{直接进入解}；用「节点值代替积分」（集中质量）会把插值误差与
积分误差混在一起，让收敛阶无法归因。

\field{关键结论} 只接受次数 $\le2$ 的多项式 $f$（有闭式），其余需要求积器（不做）。
成立条件：$f$ 多项式；单元几何规则。

\field{判据} 1D 节点精确性（F6）：$u=x(1-x)$、$f=2$ 时 P1 解在节点上与解析解\textbf{逐点相等}（机器精度）——
载荷、边界、单元刚度任一错，这条立刻红。

\field{前置} F2、F3。

\field{常见误解} 以为「一致性质量矩阵不如集中质量好」。集中质量是对角矩阵、求解快，
但它\textbf{对应另一个离散问题}；混用会让「二阶收敛」的归因失效。
\end{know}

\begin{know}{F5}{Dirichlet 消元与保持对称}{未建：math/assemble.hpp}
\field{是什么} 对边界自由度：把对应行列清零、对角置 1、右端置为边界值，并把边界值对内部的贡献从 $b$ 里扣掉。

\field{为什么需要} 不处理边界，$A$ 有非平凡零空间（常向量），CG 必失败（F9）；
处理不当，$A$ 变不对称，CG 的整个理论前提崩掉。

\field{关键结论} 「行、列都清零 + 对角置 1」保持\textbf{对称}；只清行会破坏对称性。
成立条件：Dirichlet 值已知；边界自由度集合正确（F1）。

\field{判据} 处理后 $A$ 仍对称、对角非零、$x$ 的边界分量等于给定值；未加边界时内部节点行和为 0。

\field{前置} F1、F2。

\field{常见误解} 以为「把边界节点从矩阵里删掉」更干净。删行删列会改变未知量编号（要重排索引），
工程上更容易出错；清零法保留维度、保持结构。
\end{know}

\begin{know}{F6}{收敛阶 2 与 1D 的节点精确性}{未建：math/assemble.hpp}
\field{是什么} P1 元在 $L^{2}$ 下有二阶收敛：$\norm{u-u_h}_{L^{2}}\le Ch^{2}\abs{u}_{H^{2}}$。
1D 的特殊性：节点上的值\textbf{精确}（不用减 $O(h^{2})$）。

\field{为什么需要} 它是「离散解是否可信」的唯一客观判据；1D 的节点精确性是整条链最锋利的验收工具。

\field{关键结论} 阶 2 的来源：$L^{2}$ 误差 $\sim h\cdot$（插值误差 $O(h)$ 的一阶矩），
两个 $h$ 相乘。1D 节点精确的原因：Green 函数对称性使得节点误差的主项恒为零——这是 1D 的偶然，\textbf{2D 不成立}。
成立条件：解足够光滑（$H^{2}$）；网格拟一致。

\field{判据} $n=16/32/64/128$ 三次 $h$ 减半的误差比值落在 $[3.8,4.2]$；
2D 用 $u=x(1-x)y(1-y)$、$f=2y(1-y)+2x(1-x)$，阶 2 但节点\textbf{不}精确。

\field{前置} A5、F2、F4。

\field{常见误解} 以为「1D 节点精确说明我的离散是超收敛的」。它是 1D 的结构性巧合；
照搬到 2D 会得到一个必然失败的判据。
\end{know}

\begin{know}{F7}{CSR 与 SpMV 的复杂度}{未建：linalg/sparse.hpp}
\field{是什么} CSR 用三个数组：\texttt{values}（非零元）、\texttt{colIdx}（列号）、
\texttt{rowPtr}（每行起点，长 $n+1$，\texttt{rowPtr[n]}=\text{nnz}）。

\field{为什么需要} 有限元的 $A$ 每行只有几个非零；稠密存储的 $O(n^{2})$ 与「想解多大的问题」直接冲突。

\field{关键结论} 一次 SpMV 是 $O(\text{nnz})$，而稠密是 $O(n^{2})$；对 1D 三对角 $\text{nnz}=3n-2$，
所以每次迭代 $O(n)$ 而非 $O(n^{2})$。成立条件：\texttt{rowPtr} 单调不减且首末正确；
稀疏没有「结构零」以外的假设。

\field{判据} 与稠密乘逐点对拍（随机 $100\times100$、密度 5%、$\le10^{-14}$）；
三对角手算 \texttt{rowPtr}=$[0,2,5,8,\dots]$；边界形状 $1\times1$、单行、全零行、$\text{nnz}=0$ 不许崩。

\field{前置} F1。

\field{常见误解} 以为「稀疏存储一定更快」。$n$ 小或密度高时，索引开销会反超；
判据是「实际耗时比」而不是「理论复杂度」。
\end{know}

\begin{know}{F8}{Krylov 子空间与共轭方向}{未建：linalg/cg.hpp}
\field{是什么} CG 在 Krylov 子空间 $\mathcal K_k=\mathrm{span}\{b,Ab,\dots,A^{k-1}b\}$ 中找
使 $\norm{e_k}_{A}$ 最小的 $x_k$；方向 $p_k$ 满足 $A$-共轭（$p_i^{H}Ap_j=0$，$i\ne j$）。

\field{为什么需要} 它是「迭代解」的原型：只用一个矩阵-向量乘，不存矩阵，内存 $O(n)$。

\field{关键结论} 三项递推：$r_k=b-Ax_k$、$\alpha_k=\frac{r_k^{H}r_k}{p_k^{H}Ap_k}$、
$x_{k+1}=x_k+\alpha_kp_k$、$p_{k+1}=r_{k+1}+\beta_kp_k$。共轭方向保证「每个方向只走一次」，
故理论上 $n$ 步精确收敛。成立条件：$A$ 对称正定（否则 $\norm{\cdot}_{A}$ 不是范数）。

\field{判据} 与直接解对拍，残差 $\le10^{-10}$；迭代数只记录不设阈值（1D 无预条件约 $O(n)$ 量级）；
把 $\alpha$ 写成反号必须立刻发散——这是注入对照。

\field{前置} F7、L4、N8。

\field{常见误解} 以为「共轭梯度就是最速下降的改进版」。最速下降每步正交于残差、会「锯齿」；
CG 每步 $A$-正交于所有历史方向，所以能在一个 $n$ 维空间里走 $n$ 步精确到达。
\end{know}

\begin{know}{F9}{CG 的收敛率与「为什么必须 SPD」}{未建：linalg/cg.hpp}
\field{是什么} 误差界：
\[
\norm{e_k}_{A}\le2\left(\frac{\sqrt\kappa-1}{\sqrt\kappa+1}\right)^{k}\norm{e_0}_{A},
\]
$\kappa=\lambda_{\max}/\lambda_{\min}$。

\field{为什么需要} 它决定「迭代多少次」，也解释「为什么预处理是必要的」（把 $\kappa$ 压小）。

\field{关键结论} $\kappa$ 大时收敛慢（$\kappa=10^{4}$ 需要约 $100$ 次以上量级）；
非 SPD 时 $p^{H}Ap$ 可能为 0 或负，$\alpha$ 无意义 → 残差停滞或上升。
成立条件：$A$ 对称正定；精确算术；浮点上收敛会因舍入而偏离界。

\field{判据} 非 SPD 的 \texttt{apply}（例如未加边界的纯拉普拉斯）必须表现为残差停滞/上升并被
最大迭代数截断，\textbf{不许死循环}。

\field{前置} F8、A4、N8。

\field{常见误解} 以为「CG 对任何矩阵都能用，只是慢一点」。对非对称矩阵它连「下降方向」都无法保证，
会直接崩；这就是为什么对称性（F5）是硬约束。
\end{know}
